\section{The argument and its sources}\label{sec:orientation}

Write \(\xi=\zeta(5)=\sum_{j\geq1}j^{-5}\).
The useful starting point in Persiflage's account is a positive measure
whose polynomial moments are rational, while certain values of its
Stieltjes transform contain \(\xi\). Multiplying the measure by a
particular rational function then produces a matrix with two simultaneous
descriptions: its entries are affine expressions in \(\xi\), and it is
a Gram matrix. We follow that organizing perspective
\cite[The formulas of Euler and Hermite; Hankel Determinants]{persiflage}.

The historical distinction matters. Pr\'evost's Theorems 1--2
\cite{prevost} supply the underlying transform and moments; his earlier
work on Pad\'e approximation, and subsequent work with Rivoal
\cite{prevost1996,prevost-rivoal}, explain their role in rational
approximation. Rationally modifying an orthogonality measure belongs to
the theory developed by Uvarov, with a determinant formula proved by
Krattenthaler \cite{uvarov,krattenthaler}. Zudilin's determinantal irrationality
criterion \cite{zudilin} is a direct methodological antecedent.
Brown's integral pairings and coefficient-normalization framework
\cite[Sections 3 and 7.4--7.5]{brown} provide further context.
We do not need these general determinant formulas as unproved steps:
the particular identities used here are established below.

The extra work for this construction is quantitative. The numerator
of the rational modifier improves divisibility, but also changes the
real integral. Both effects must be estimated for the same determinant.
Fauzan makes this choice explicit \cite[Sections 2--6]{fauzan}.
For the local estimates we use Firsching's revised finite
partial-fraction argument \cite{firsching}. In particular, the correction
of bounded rank is charged its full rank, and the prime cutoff is 400.
No inequality with the smaller, disputed rank loss is needed.

\subsection{What the estimates must accomplish}

The desired polynomials have integer coefficients and satisfy
\begin{equation}\label{eq:goal}
 \deg Q_n\leq37n,\qquad
 0<Q_n(\xi)\leq e^{-cn^2}
\end{equation}
for a fixed \(c>0\) and all sufficiently large \(n\).
This form of the final contradiction is used by Fauzan
\cite[Section 1.1]{fauzan}; compare Zudilin's determinant argument
\cite[Section 2]{zudilin}.

Suppose \(\xi=a/b\), where \(a\in\Z\) and \(b\in\Z_{>0}\).
Writing \(Q_n(X)=\sum_{j=0}^{37n}q_jX^j\), with zero coefficients added
if necessary, gives
\[
 b^{37n}Q_n(a/b)=\sum_{j=0}^{37n}q_j a^j b^{37n-j}\in\Z.
\]
It is positive by \eqref{eq:goal}, whereas
\[
 0<b^{37n}Q_n(a/b)
 \leq\exp(37n\log b-cn^2)\longrightarrow0.
\]
The denominator \(b\) is fixed. A positive integer cannot tend to zero.
This proves irrationality once \eqref{eq:goal} has been established.

We construct a rational polynomial \(F_K\), then a positive rational
multiplier \(m_K\) such that \(Q_n=m_KF_K\) for \(K=40n\).
Positivity will follow before any size estimate is used.
The analytic calculation bounds \(\log F_K(\xi)\) by a negative
constant times \(K^2\), up to lower-order terms.
The arithmetic calculation bounds \(\log m_K\) by a smaller positive
constant times \(K^2\). Their sum is negative.

\subsection{Where the details live}

The integral representation is derived in full. Sections 5--7 then keep
the main inequalities and the reasons behind them together, so that
the real decay can be compared directly with the arithmetic cost.
The supporting calculations are in the separate
\emph{detailed companion}, supplied as
\texttt{zeta5\_detailed\_companion.pdf}: its Section 5 proves the analytic
bounds, Sections 6--7 prove the local and global arithmetic bounds,
and Appendix C gives the exact arithmetic table.
Those references supply proofs of the stated estimates, not additional
assumptions. The unchanged finite certificates and source ledger accompany
both versions. No formal verification of this prose or local rebuild of
the complete Lean development is claimed.
